\documentclass[10pt,a4paper]{amsart}
\usepackage[margin=19mm]{geometry}
\usepackage{amsmath,amssymb,mathtools}
\usepackage[scr=rsfs,cal=euler]{mathalfa}
\usepackage{xfrac}
\usepackage[unicode,hidelinks,bookmarks=false]{hyperref}
\newcommand{\ie}{i.e.\ }
\newcommand{\cf}{cf.\ }
\newcommand{\resp}{respectively\ }
\newcommand{\Subsection}{\subsection}
\providecommand{\nobreakdash}{\nobreak\hbox{-}}
\setlength{\emergencystretch}{2em}
\multlinegap0pt
\usepackage{bm}
\usepackage{bbm}
\usepackage{dsfont}
\usepackage{xspace}
\usepackage{cancel}
\usepackage{mathtools}
\usepackage{amsthm}
\makeatletter
\@ifundefined{thm}{\newtheorem{thm}{Theorem}}{}
\@ifundefined{prop}{\newtheorem{prop}[thm]{Proposition}}{}
\makeatother
\makeatletter
\DeclareMathOperator{\N}{N}
\newcommand{\cT}{\mathcal{T}}\newcommand{\cU}{\mathcal{U}}
\newcommand{\convp}{\stackrel{P}{\longrightarrow}}
\expandafter\ifx\csname enumerateA\endcsname\relax
\newenvironment{enumerateA}{\begin{enumerate}[\upshape (A)]\setlength{\itemsep}{1ex}}{\end{enumerate}}
\fi
\expandafter\ifx\csname enumeratea\endcsname\relax
\newenvironment{enumeratea}{\begin{enumerate}[\upshape a)]\setlength{\itemsep}{1ex}}{\end{enumerate}}
\fi
\expandafter\ifx\csname enumeratei\endcsname\relax
\newenvironment{enumeratei}{\begin{enumerate}[\upshape i)]}{\end{enumerate}}
\fi
\expandafter\ifx\csname enumeraten\endcsname\relax
\newenvironment{enumeraten}{\begin{enumerate}[\upshape 1.]\setlength{\itemsep}{1ex}}{\end{enumerate}}
\fi
\renewcommand{\ge}{\geqslant} 
\expandafter\ifx\csname inparaenuma\endcsname\relax
\newenvironment{inparaenuma}{\begin{inparaenum}[\upshape \bfseries (a) ]}{\end{inparaenum}}
\fi
\expandafter\ifx\csname inparaenumn\endcsname\relax
\newenvironment{inparaenumn}{\begin{inparaenum}[\upshape \bfseries (1) ]}{\end{inparaenum}}
\fi
\renewcommand{\le}{\leqslant} 
\DeclareMathOperator{\pr}{\mathds{P}}
\makeatother
\title[Evolution of recursive trees with limited memory]{Evolution of recursive trees with limited memory}
\author{Omer Angel, Shankar Bhamidi, Serte Donderwinkel, Neeladri Maitra, Akshay Sakanaveeti}
\date{Source: arXiv:2510.18856v1 (CC BY 4.0)}
\begin{document}
\maketitle

\noindent\emph{Source and licence.} Evolution of recursive trees with limited memory. Version arXiv:2510.18856v1 (CC BY 4.0), distributed under the Creative Commons Attribution 4.0 International licence, as recorded in the arXiv licence metadata for this exact version and shown on the abstract page https://arxiv.org/abs/2510.18856v1. Full rights notice: licenses/arxiv-2510.18856-CC-BY-4.0.txt.\par\smallskip

\noindent\textit{Editorial scope.} Counted unit: the Prop whose source label is \texttt{prop:dist\_limit} in the article (source0.tex lines 1282--1285, source label \texttt{prop:dist\_limit}), delivered together with the article's own complete original proof of it (source0.tex lines 1408--1426, one \texttt{proof} environment of 19 lines). No mathematics is written, repaired or re-derived: statement and proof are the article's own text line for line. The proof is detached from the statement in the source: the article states the result at lines 1282--1285 and prints its proof at lines 1408--1426, and the proof environment's own header names the statement, so the association is the article's own. Required context retained uncounted: thm:ghp\_meso (lines 688--693). Every definition, display and label of those lines is retained unchanged. Omitted from this package and not redistributed: 1--687, 694--1281, 1286--1407, 1427--1578. Nothing inside a retained excerpt is omitted or altered: every retained body is delivered exactly as the article's lines are written, so no omission receipt of any kind accompanies this package. Citations: the retained excerpts carry no citation command at all, so no bibliography entry of the article is transcribed and no reference is redistributed. Reference adaptation: none was needed and none was made; every reference inside the delivered unit resolves inside it, so no unresolved reference or citation is produced. Printed slips kept literal: no printed word, symbol, hypothesis or number of the counted statement or of its proof was altered. Layout and class: the article's own class is \texttt{amsart} with packages including geometry, authblk, comment, systeme, mathrsfs, xfrac, hyperref, amsmath, amssymb, amsthm, amsfonts, amsbsy, latexsym, dsfont; the wrapper is the lane's pinned \texttt{amsart} editorial wrapper at 10pt on A4, which restates the theorem-like environments the retained text needs and adds the article's own macro definitions that the retained text needs (\texttt{\textbackslash N}, \texttt{\textbackslash cT}, \texttt{\textbackslash convp}, \texttt{\textbackslash ge}, \texttt{\textbackslash le}, \texttt{\textbackslash pr}), with extra wrapper packages (bm, bbm, dsfont, xspace, cancel, mathtools). The delivered numbering is the unit's own. Assets: no retained span contains a figure environment, a graphics-inclusion command, a PDF-page inclusion, a table or a TikZ picture; no image file is copied into this package, the package declares no asset dependency and no image file is redistributed with it. Licence: version arXiv:2510.18856v1 of this article is distributed under the Creative Commons Attribution 4.0 International licence, as recorded in the arXiv licence metadata for this exact version and shown on the abstract page https://arxiv.org/abs/2510.18856v1. A full rights notice is delivered at licenses/arxiv-2510.18856-CC-BY-4.0.txt. Characters: every retained excerpt is pure ASCII, so the delivered engine, XeLaTeX with its default Latin Modern text font, reports no missing character.
\begin{thm}\label{thm:ghp_meso}
In the mesoscopic regime when $\beta\in (0,1/2)$ we have 

\[\frac{1-\beta}{2} n^{-(1-\beta)}\mathcal{T}_n \convp L  \]
in the Gromov--Hausdorff topology, with $L$ a line segment of length $1$. 
\end{thm}

\begin{prop}\label{prop:dist_limit}
For $\beta \in (0,1)$, it holds that 
\[n^{-(1-\beta)} d(1,n)\convp \frac{2}{1-\beta}. \]
\end{prop}

\begin{proof}[Proof of Theorem \ref{thm:ghp_meso}]
Let $\beta<1/2$. In Proposition \ref{prop:dist_limit} we showed that $n^{-(1-\beta)}d_{\mathcal{T}_n}(1,n)\convp \frac{2}{1-\beta}$ as $n\to\infty$.
Now, fix $\varepsilon>0$ small enough so that $\beta+\varepsilon<1-\beta$. We will show that with high probability, for each $m<n$, the distance from vertex $m$ to the ancestral line of $n$ satisfies $d(m,m\wedge n)<n^{\beta+\varepsilon}$, which shows that, with high probability, the path from the root to vertex $n$ is at Hausdorff distance $o(n^{1-\beta})$ from $\cT_n$ and thus proves the theorem. 

Let $G_n\sim \operatorname{Geom}(n^{-\beta})$ a Geometric random variable with for $k\in \N$, $\pr(G_n>k)=(1-n^{-\beta})^k$. We will show that for each $m<n$ \begin{equation}\label{eq:geom_stochdom}d(m,m\wedge n)\prec_{st}G_n+1.\end{equation} 
Then, 
\[\pr\left(\exists m<n: d(m,m\wedge n)>n^{\beta+\varepsilon} \right) \le n (1-n^{-\beta})^{n^{\beta+\epsilon}}\le n\exp(-n^\epsilon)\to 0\]
as $n\to \infty$. Thus we focus on showing \eqref{eq:geom_stochdom}. 

We condition on the unique path
\[P_n=(1=L_n(d(1,n))<L_n(d(1,n)-1)<\dots< L_n(0)=n)\] in $\mathcal{T}_n$ from vertex $1$ to vertex $n$. We will show \eqref{eq:geom_stochdom} by iteratively revealing the ancestors of vertex $m$ and showing that, at each step, the next ancestor is on $P_n$ with probability at least $n^{-\beta}$. Formally, we will bound
\begin{equation}\label{eq:prob_ancestor} \pr(L_m(k+1)\in P_n | P_n, L_m(1), \dots, L_m(k))\end{equation}
from below by $n^{-\beta}$.

We use that for any $2\le m\le n$, $[j(m-1), m-1]\cap P_n$ is non-empty. Indeed, let $\ell$ be maximal such that $L_n(\ell)\ge m$. If $L_n(\ell)=m$, then the claim follows immediately. Otherwise, \[L_n(\ell+1)\in [j(L_n(\ell)-1) ,  m-1]\subset [j(m-1), m-1]  \] because $j(n)$ is non-decreasing. 

Then, to bound \eqref{eq:prob_ancestor} we see that either $L_m(k)\in P_n$, in which case \eqref{eq:prob_ancestor} equals $1$. Otherwise, on the conditioning, $L_m(k+1)$ is a unique element from $[j(L_m(k)-1), L_m(k)-1]$. By our observation, at least one element of this set is also in $P_n$. Moreover, the set has size at most $n^{\beta}$, so \eqref{eq:prob_ancestor} is at least $n^{-\beta}$. This implies \eqref{eq:geom_stochdom}. 

\end{proof}

\end{document}
